\chapter{Marco Teórico y Metodología}
\label{cap:MarcoTeorico}

% \epigraph{La perfección no se alcanza cuando no hay nada más que añadir, sino cuando no hay nada más que quitar.}{Antoine de Saint-Exupéry}

\lettrine[lraise=-0.1, lines=2, loversize=0.2]{E}{n} este capítulo se establece el marco teórico y metodológico sobre el que se articula el sistema propuesto. En primer lugar, se formula el modelo matemático del canal de radiofrecuencia y se caracteriza el régimen de baja relación señal a ruido que constituye el principal desafío del problema. A continuación, se presenta el conjunto de datos seleccionado como eje experimental, \textit{NoisyUAV v2}, y se analiza en detalle su estructura y propiedades espectrales. Seguidamente, se justifica el procesamiento directo de las secuencias I/Q frente a las representaciones tiempo-frecuencia clásicas, se formaliza el problema del ruido de etiquetas inherente al formato de los datos y, finalmente, se describe el detector analítico de ráfagas basado en entropía espectral de Shannon que actúa como preprocesador del clasificador neuronal.

\section{Modelo del Sistema y Formulación del Canal RF}
El diseño de un sistema robusto de detección pasiva requiere, en primer lugar, una caracterización rigurosa del entorno electromagnético y de las señales involucradas. En este proyecto, se asume un escenario de monitorización donde un receptor de radio definida por \textit{software} (\gls{sdr}) escanea de forma continua la banda Industrial, Científica y Médica (\gls{ism}) de 2.4\,GHz.

La señal captada por la antena receptora en el instante de tiempo $t$, denotada como $y(t)$, puede modelarse matemáticamente como la superposición de la emisión del objetivo, las alteraciones del canal y las emisiones espurias del entorno \cite{rout_novel_2025, ezuma_detection_2020}. En su representación de banda base equivalente, el modelo del sistema se define mediante la ecuación:

\begin{equation}
    y(t) = \big(h * x\big)(t) + i(t) + n(t)
    \label{eq:modelo_sistema}
\end{equation}

Donde cada término representa una magnitud física discreta del canal inalámbrico:
\begin{itemize}
    \item $x(t)$ es la señal compleja en banda base transmitida por el UAV o su radiocontrol. Típicamente, estas señales emplean esquemas de modulación de espectro ensanchado por salto en frecuencia (\textit{Frequency Hopping Spread Spectrum}, FHSS) para mitigar interferencias y evitar intercepciones.
    \item $h(t)$ representa la respuesta al impulso compleja del canal físico en banda base equivalente. Este término compacta los fenómenos de atenuación por propagación (\textit{path loss}), rotación fasorial de fase y desvanecimiento destructivo/constructivo inter-símbolo, modelando las transformaciones que sufre la envolvente compleja de la señal antes de la etapa de digitalización en el receptor.
    \item $i(t)$ representa la interferencia co-canal. Este término agrupa las emisiones de otros dispositivos coexistentes en la banda de interés, principalmente ráfagas de redes IEEE 802.11 (Wi-Fi) y Bluetooth \cite{ezuma_detection_2020}.
    \item $n(t)$ modela el ruido térmico inherente a la electrónica del receptor, habitualmente caracterizado como ruido blanco gaussiano aditivo (\textit{Aditive White Gaussian Noise}, AWGN) con media nula y varianza $\sigma_n^2$ \cite{lacy_machine_2024}.
\end{itemize}

Para su procesamiento digital en las arquitecturas de aprendizaje profundo, la señal analógica $y(t)$ es muestreada a una frecuencia $F_s$, obteniendo la secuencia discreta $y[n]$. Para preservar la coherencia de fase espacial y aislar las firmas electromagnéticas del transmisor (\textit{hardware impairments}), la señal se procesa directamente como un vector de valores complejos compuesto por sus componentes en fase ($I$) y en cuadratura ($Q$) \cite{valkama_blind_2005, zheng_deep_2025}:
\begin{equation}
    y[n] = y_I[n] + j y_Q[n]
    \label{eq:iq_discreta}
\end{equation}

\subsection{Definición del régimen de baja Relación Señal a Ruido (SNR)}
El desafío fundamental abordado en este proyecto radica en la detección de la ráfaga $x(t)$ cuando su contribución energética es eclipsada por los términos $i(t)$ y $n(t)$. La Relación Señal a Ruido e Interferencia (Signal-to-Interference-plus-Noise Ratio, SINR), o de forma generalizada SNR en el contexto de conjuntos de datos contaminados, se define paramétricamente como:
\begin{equation}
    \text{SNR (dB)} = 10 \log_{10} \left( \frac{P_x}{P_i + P_n} \right)
    \label{eq:snr_calculo}
\end{equation}

Donde $P_x$, $P_i$ y $P_n$ denotan las potencias medias de la señal de interés, la interferencia y el ruido térmico, respectivamente. Los detectores de energía clásicos y las arquitecturas de aprendizaje automático convencionales sufren una degradación de rendimiento severa cuando la \gls{snr} desciende por debajo de los 0\,dB, ya que la envolvente de la señal $y[n]$ se vuelve estadísticamente indistinguible de la distribución gaussiana del ruido \cite{lacy_machine_2024}.

Como se detallará en las secciones posteriores, el objetivo de la metodología propuesta es extraer características no lineales directamente de los tensores I/Q descritos en la Ecuación \ref{eq:iq_discreta} para garantizar la detección en regímenes críticos donde $\text{SNR} < -6$\,dB.

\section{El Conjunto de Datos NoisyUAV v2}
\label{sec:dataset_noisyuav}

Tal y como se concluye del análisis crítico de los repositorios existentes presentado en el Capítulo~\ref{cap:estado_arte}, ninguno de ellos satisface los requisitos necesarios para abordar la detección en el límite del ruido térmico. Por ello, se ha seleccionado el \textit{dataset} NoisyUAV v2 como el eje experimental de este proyecto. Esta base de datos se diseñó con el objetivo explícito de evaluar la robustez algorítmica en entornos de muy baja SNR y con interferencias reales.

Autores y artículo: Glüge et al. (2024) \cite{gluge_robust_2024}.


Arquitectura de adquisición de RF: captación realizada empleando una plataforma SDR basada en componentes comerciales (COTS, \textit{Commercial Off-The-Shelf}), concretamente el USRP B205mini-i. Se operó a una frecuencia de muestreo ($F_s$) de 14\,MHz en la banda de 2.4\,GHz.

Volumen, topología y formato de los datos: compone 17\,.744 grabaciones almacenadas como vectores analíticos I/Q. Cada archivo representa una ventana de observación continua de 75\,ms (1.048.576 muestras). La taxonomía abarca seis modelos bajo distintos protocolos FHSS (DJI, Futaba, Graupner, Taranis, Turnigy) y una clase exclusiva de ruido empírico (Wi-Fi, Bluetooth, AWGN).

Análisis de la hostilidad electromagnética: las muestras se estratifican en tres regímenes:
\begin{itemize}
    \item Grupo A (SNR $\ge$ 10\,dB): potencia significativamente superior al suelo de ruido. Sin embargo, las interferencias presentan perfiles energéticos equiparables a las transmisiones, incrementando la ambigüedad.
    \item Grupo B ($-6$\,dB $\le$ SNR $<$ 10\,dB): complejidad moderada, donde la coexistencia co-canal degrada la integridad fasorial.
    \item Grupo C (SNR $<$ $-6$\,dB): entornos hostiles críticos, con ráfagas capturadas hasta los $-20$\,dB. Las señales quedan completamente enmascaradas bajo el ruido térmico.
\end{itemize}

Justificación de la selección: constituye el único conjunto de datos que satisface íntegramente los requisitos definidos. Los 14\,MHz de $F_s$ evitan la redundancia computacional, los 75\,ms continuos permiten el escaneo activo, y la presencia física de regímenes de SNR negativa garantiza una validación robusta de las arquitecturas propuestas.

\subsubsection{Inspección del conjunto de datos seleccionado}

Para comprender el escenario de detección, resulta fundamental analizar la naturaleza fasorial y espectral de las señales. Como detalla la Tabla \ref{tab:distribucion_muestras}, la distribución equitativa interclase permite un entrenamiento robusto.

\begin{table}[htbp]
\centering
\small
\caption{Distribución de muestras en el conjunto de datos NoisyUAV v2.}
\label{tab:distribucion_muestras}
\resizebox{\linewidth}{!}{%
\begin{tabular}{lccccc}
\toprule
\textbf{Origen} & \textbf{ID clase} & \textbf{Muestras} & \textbf{$f_c$ (GHz)} & \textbf{$\Delta f$ (MHz)} & \textbf{Duración (ms)} \\ \midrule
Dron DJI Phantom GL300F Pro & 0 & 1280 & 2.44175 & 1.2 & 2.18\\
Dron Futaba T7C & 1 & 3472 & 2.44175 & 2.0 & 1.70\\
Dron Futaba T14SG & 2 & 801 & 2.44175 & 3.1 & 1.40\\
Graupner mx-16 & 3 & 801 & 2.44175 & 1.0 & 1.9/3.7\\
Ruido & 4 & 8872 & - & - & - \\
Dron Taranis ACCST & 5 & 1663 & 2.44000 & 1.5 & 3.1/4.4\\
Dron Turnigy 9X & 6 & 855 & 2.44500 & 2.0 & 1.30\\ \bottomrule
\textbf{Total} & & \textbf{17.744} & & & \\ \bottomrule
\end{tabular}
}
\end{table}

1. Variabilidad interclase y regímenes de operación: un sistema robusto debe discriminar protocolos FHSS independientemente de la potencia recibida. La Figura \ref{fig:comparativa_clases_completa} ilustra el desvanecimiento morfológico de las ráfagas a medida que la SNR decae, dificultando la extracción de características basadas en umbrales de energía.

\enlargethispage{1.5cm} % TRUCO 1: Alargamos el folio virtualmente 2 cm para que quepa la figura ampliada

\begin{figure}[!htbp]
    \centering
    \begin{subfigure}[b]{0.49\linewidth} % TRUCO 2: Aumentamos el ancho
        \centering
        \includegraphics[width=\linewidth]{figuras/seccion3/comparacion_clase_28dB.png}
        \caption{Régimen favorable (28\,dB).}
    \end{subfigure}
    \hfill
    \begin{subfigure}[b]{0.49\linewidth}
        \centering
        \includegraphics[width=\linewidth]{figuras/seccion3/comparacion_clase_10dB.png}
        \caption{Régimen intermedio (10\,dB).}
    \end{subfigure}

    \vspace{0.3cm} % TRUCO 3: Espacio negativo para que al crecer a lo ancho no se separen tanto a lo alto

    \begin{subfigure}[b]{0.49\linewidth}
        \centering
        \includegraphics[width=\linewidth]{figuras/seccion3/comparacion_clase_0dB.png}
        \caption{Régimen moderado (0\,dB).}
    \end{subfigure}
    \hfill
    \begin{subfigure}[b]{0.49\linewidth}
        \centering
        \includegraphics[width=\linewidth]{figuras/seccion3/comparacion_clase_-12dB.png}
        \caption{Régimen crítico ($-12$\,dB).}
    \end{subfigure}
    
    % \vspace{-0.2cm} % Pegamos el texto del pie a las figuras
    \caption{Comparación por clase en función de la degradación progresiva de la relación señal a ruido.}
    \label{fig:comparativa_clases_completa}
\end{figure}

% \FloatBarrier % Tu barrera se mantiene intacta

2. Análisis multidominio de la degradación: para ilustrar el impacto físico del ruido térmico, la Figura \ref{fig:degradacion_target0_total} expone la evolución del \textit{Target 0} en los dominios fasorial (I/Q), temporal (envolvente) y tiempo-frecuencia (espectrograma y PSD). Esta transición evidencia la paulatina corrupción de la señal a medida que la SNR baja.

% \enlargethispage{2cm} % Damos margen vertical extra por si al ensanchar crecen mucho a lo alto

\begin{figure}[!htbp]
    \centering
    \begin{subfigure}[b]{0.495\linewidth} % Apuramos al máximo el ancho de la columna
        \centering
        \includegraphics[width=\linewidth]{figuras/seccion3/ej_dataset_target0_30db.png}
        \caption{\textit{Target 0} a 30\,dB.}
    \end{subfigure}
    \hfill
    \begin{subfigure}[b]{0.495\linewidth}
        \centering
        \includegraphics[width=\linewidth]{figuras/seccion3/ej_dataset_target0_0db.png}
        \caption{\textit{Target 0} a 0\,dB.}
    \end{subfigure}

    \vspace{0.3cm} % Reducimos un poco el espacio central

    \begin{subfigure}[b]{0.495\linewidth}
        \centering
        \includegraphics[width=\linewidth]{figuras/seccion3/ej_dataset_target0_-12db.png}
        \caption{\textit{Target 0} a $-12$\,dB.} % He corregido el signo menos matemático aquí
    \end{subfigure}
    \hfill
    \begin{subfigure}[b]{0.495\linewidth}
        \centering
        \includegraphics[width=\linewidth]{figuras/seccion3/ej_dataset_target0_-20db.png}
        \caption{\textit{Target 0} a $-20$\,dB.} % He ajustado el texto al nombre del archivo (-20dB)
    \end{subfigure}
    \caption{Evolución multidominio de la clase objetivo 0 en función del nivel de ruido del canal.}
    \label{fig:degradacion_target0_total}
\end{figure}

% \clearpage % Este comando es un muro absoluto: empuja todo lo que venga debajo a la página siguiente

3. Caracterización de las interferencias y falsos positivos: la Figura \ref{fig:degradacion_ruido_total} detalla el comportamiento espectral de la clase correspondiente a ruido e interferencia. Se observa cómo, a niveles altos de SNR, las ráfagas espurias Wi-Fi presentan firmas energéticas que son susceptibles de provocar falsos positivos en algoritmos clásicos de detección pasiva.

\begin{figure}[!htbp]
    \centering
    \begin{subfigure}[b]{0.495\linewidth}
        \centering
        \includegraphics[width=\linewidth]{figuras/seccion3/ej_dataset_target4_30db.png}
        \caption{Clase ruido a 30\,dB.}
    \end{subfigure}
    \hfill
    \begin{subfigure}[b]{0.495\linewidth}
        \centering
        \includegraphics[width=\linewidth]{figuras/seccion3/ej_dataset_target4_12db.png}
        \caption{Clase ruido a 12\,dB.}
    \end{subfigure}

    \vspace{0.3cm} % Espacio reducido para compensar el aumento de altura

    \begin{subfigure}[b]{0.495\linewidth}
        \centering
        \includegraphics[width=\linewidth]{figuras/seccion3/ej_dataset_target4_0db.png}
        \caption{Clase ruido a 0\,dB.}
    \end{subfigure}
    \hfill
    \begin{subfigure}[b]{0.495\linewidth}
        \centering
        \includegraphics[width=\linewidth]{figuras/seccion3/ej_dataset_target4_-12db.png}
        \caption{Clase ruido a $-12$\,dB.}
    \end{subfigure}
    \caption{Análisis de las emisiones interferentes y ruido de fondo a distintos regímenes de potencia.}
    \label{fig:degradacion_ruido_total}
\end{figure}

\section{Justificación del Procesamiento Directo I/Q frente a Representaciones de Tiempo-Frecuencia}
\label{sec:justificacion_procesamiento_IQ}

El paradigma predominante en la literatura reciente orientada a la clasificación de señales electromagnéticas mediante aprendizaje profundo, tal y como se detalla en el Capítulo \ref{cap:estado_arte}, se fundamenta en la proyección de la serie temporal unidimensional hacia una representación bidimensional en el dominio tiempo-frecuencia. En este contexto, el estimador de densidad espectral más empleado es el espectrograma \cite{gluge_robust_2024, zhao_drone_2024}. Matemáticamente, esta matriz se computa evaluando la magnitud al cuadrado de la Transformada de Fourier a Corto Plazo (\textit{Short-Time Fourier Transform}, STFT), denotada por la expresión:

\begin{equation}
    S[m,k] = \left| \sum_{n=0}^{N_w-1} y[n+mH] w[n] e^{-j\frac{2\pi}{N_w}kn} \right|^2
    \label{eq:espectrograma}
\end{equation}

En esta formulación, $w[n]$ representa la función de la ventana aplicada para mitigar la fuga espectral, $N_w$ determina la longitud de dicha ventana de análisis, y el parámetro $H$ rige el desplazamiento discreto entre segmentos temporales consecutivos. Este último factor controla de manera unívoca la tasa de solapamiento de las ventanas, dictando la resolución temporal del estimador frente a la evolución transitoria de la señal.

La motivación analítica que subyace a esta transformación espacial radica en la posibilidad de tratar la matriz espectral resultante como una imagen digital fotométrica. Este isomorfismo permite importar e implementar arquitecturas de redes neuronales convolucionales que fueron concebidas para tareas de reconocimiento de imágenes. En dicho marco teórico, el propósito de las capas de convolución es recorrer la matriz espacial bidimensional para identificar jerarquías de características topológicas, asimilando las variaciones de densidad espectral a gradientes espaciales, bordes o patrones de textura visual. 

No obstante, la asimilación conceptual de una señal de radiofrecuencia a un mapa de píxeles carece de rigor físico, induciendo deficiencias estructurales en el modelo. Concretamente, la adopción del espectrograma como bloque de entrada impone tres limitaciones fundamentales que comprometen severamente la detección en regímenes donde la relación señal a ruido desciende por debajo de los 0\,dB:

\begin{itemize}
    \item Destrucción de la coherencia de fase: la operación del valor absoluto en la Ecuación \ref{eq:espectrograma} descarta la fase instantánea de la señal. En regímenes críticos ($\text{SNR} \le -6$\,dB), las transiciones de fase inherentes a la modulación y las firmas de \textit{hardware} de los osciladores locales (\textit{RF fingerprinting}) constituyen discriminantes estadísticamente vitales frente al ruido gaussiano que resultan irrecuperables tras la transformación espectral \cite{zheng_deep_2025}.
    \item Límite de Heisenberg-Gabor: la STFT está sujeta al principio de incertidumbre tiempo-frecuencia, que establece un compromiso ineludible: una alta resolución frecuencial (ventanas largas) diluye los bordes transitorios de las ráfagas FHSS, mientras que una alta resolución temporal (ventanas cortas) ensancha artificialmente el espectro, incrementando el solapamiento con interferencias co-canal \cite{oppenheim_discrete_2014}.
    \item Complejidad espacial y carga computacional: la proyección de una serie temporal unidimensional hacia una cuadrícula bidimensional tiempo-frecuencia supone un aumento directo de la cantidad de parámetros entrenables en los modelos \cite{alam_rf-enabled_2023}. En una arquitectura CNN 2D convencional, el número de operaciones de multiplicación-acumulación (MAC) escala de forma cuadrática con el tamaño del núcleo o \textit{kernel} convolucional ($K \times K$). Para preservar la resolución de una señal muestreada a 14\,MHz (como es el caso de las señales del conjunto de datos seleccionado), los espectrogramas resultantes exigen tensores de entrada de dimensiones muy altas. En contraste, el procesamiento de la secuencia en banda base permite el uso de convoluciones unidimensionales, cuya complejidad lineal ($\mathcal{O}(N \cdot K)$) reduce drásticamente el consumo de memoria y los tiempos de inferencia \cite{zheng_deep_2025}, habilitando el despliegue del modelo en equipos alojados en el borde sin sacrificar capacidad de representación \cite{cai_toward_2024, zhang_edge_2023}.
\end{itemize}

Por consiguiente, la metodología de este proyecto descarta las representaciones espectrales destructivas y aboga por el procesamiento directo de las secuencias analíticas I/Q. Al inyectar el tensor complejo $y[n] = y_I[n] + j y_Q[n]$ directamente en la red neuronal, se delega en los núcleos convolucionales la tarea de aprender una base de filtrado adaptativa que supera los límites matemáticos de la \gls{stft} y preserva la integridad fasorial.

\section{Impacto de la Dispersión Temporal en la Corrupción de Etiquetas (\textit{Label Noise})}
\label{sec:label_noise}

La naturaleza no licenciada de la banda ISM de 2.4\,GHz induce un entorno electromagnético regido por la coexistencia estocástica de múltiples emisores asíncronos, predominantemente ráfagas de banda ancha (e.g., IEEE 802.11) y de banda estrecha (e.g., IEEE 802.15.1) \cite{ezuma_detection_2020}. En el paradigma de adquisición del conjunto de datos \textit{NoisyUAV v2}, cada registro fasorial abarca una ventana de observación macroscópica de $T_{\text{obs}} = 75$\,ms. 

La inyección directa de secuencias de 75\,ms (1.048.576 muestras) supone una carga computacional muy alta y agrava la ceguera del modelo frente a señales débiles. Esto se debe a que la información útil representa una fracción minúscula del tensor, el cual queda invariablemente dominado por el ruido térmico y las interferencias co-canal.

Para cuantificar matemáticamente este fenómeno, resulta imperativo definir el ciclo de trabajo empírico de los protocolos FHSS dentro de la ventana de observación $T_{\text{obs}} = 75$\,ms. La ocupación espacial máxima de la señal objetivo dentro del vector analítico se calcula como el cociente porcentual entre el tiempo activo total y la duración de la ventana, de acuerdo con la siguiente expresión:

\begin{equation}
    \text{Ocupación máxima (\%)} = \frac{N_{\text{burst}} \cdot \max(T_{\text{burst}})}{T_{\text{obs}}} \times 100
    \label{eq:ocupacion_maxima}
\end{equation}

Donde $N_{\text{burst}}$ representa el número de ráfagas estimadas presentes en una captura continua para un controlador específico, y $\max(T_{\text{burst}})$ denota la duración temporal de dichas ráfagas. Aplicando la Ecuación \ref{eq:ocupacion_maxima} a los datos paramétricos de la Tabla \ref{tab:distribucion_muestras} y a la inspección física de las tasas de salto de frecuencia reales, se obtienen las métricas de ocupación temporal para cada modelo documentado en el conjunto de datos, tal y como se detalla en la Tabla \ref{tab:ocupacion_temporal}.

\begin{table}[htbp]
\centering
\small
\caption{Ciclo de trabajo empírico y ocupación temporal máxima por modelo en una ventana de observación de 75\,ms.}
\label{tab:ocupacion_temporal}
\resizebox{0.9\linewidth}{!}{%
\begin{tabular}{lccc}
\toprule
\textbf{Dron / Controlador} & \textbf{Ráfagas estimadas ($N_{\text{burst}}$)} & \textbf{Duración máxima $\max(T_{\text{burst}})$ (ms)} & \textbf{Ocupación máxima (\%)} \\ \midrule
Turnigy 9X & 6 & 1.30 & 10.40\% \\
Taranis ACCST & 2 & 4.40 & 11.73\% \\
Graupner mx-16 & 2 & 3.70 & 9.86\% \\
Futaba T14SG & 2 & 3.10 & 8.26\% \\
DJI Phantom GL300F & 2 & 2.18 & 5.81\% \\
Futaba T7C & 2 & 2.00 & 5.33\% \\ \bottomrule
\end{tabular}
}
\end{table}

A partir de los resultados expuestos, se observa que la cota superior de ocupación (el escenario más favorable para la detección) corresponde al modelo Taranis ACCST. Este dispositivo transmite, en media, dos ráfagas en la ventana de grabación y cuenta con transitorios de mayor extensión temporal, alcanzando una ocupación máxima del 11.73\,\%. Resulta pertinente señalar que, incluso para el modelo Turnigy 9X, el cual presenta una tasa de salto significativamente mayor (6 ráfagas por ventana), su reducida duración individual ($1.30$\,ms) sitúa su ocupación temporal en un 10.40\,\%.

En consecuencia, incluso en el escenario analítico que asume la máxima ocupación temporal posible, más del 88.27\,\% del tensor (aproximadamente $925.574$ muestras complejas) carece de información perteneciente a la clase objetivo. Este vasto subespacio temporal está constituido exclusivamente por el suelo de ruido térmico o por las ya mencionadas ráfagas espurias procedentes de redes Wi-Fi y Bluetooth coexistentes en la banda de interés.

Metodológicamente, a esta serie temporal compleja de longitud 1.048.576 muestras se le asigna una etiqueta estática global $\tilde{y} \in \mathcal{Y}$, asumiendo de manera implícita la invariancia de la clase objetivo durante todo el intervalo de integración \cite{gluge_robust_2024}. En regímenes asintóticos de alta relación señal a ruido ($\text{SNR} \gg 0$\,dB), la densidad espectral de potencia de la señal modulada del UAV ($P_x$) domina el tensor de observación frente al umbral de interferencia ($P_i$) y al suelo de ruido térmico ($P_n$). En este caso, la etiqueta global $\tilde{y}$ constituye una cota inferior representativa de la variable aleatoria latente, garantizando que cualquier sistema de reconocimiento de patrones asocie de forma unívoca la topología de la señal con la firma electromagnética del dron. Sin embargo, cuando el canal opera en condiciones de degradación alta ($\text{SNR} < 0$\,dB), el límite físico de la transmisión sitúa a $P_x$ por debajo del umbral del ruido de fondo ($P_x < P_i + P_n$). Consecuentemente, el vector de observación queda dominado en el dominio de la probabilidad por la superposición de emisiones espurias y la varianza del proceso gaussiano (AWGN).

Al proyectar la totalidad de la secuencia temporal sobre la función de coste de un algoritmo de clasificación bajo la premisa de que toda la ventana de 75\,ms se corresponde con la clase del UAV, se induce una corrupción sistemática en el proceso de optimización. Este fenómeno, formalizado en la teoría del aprendizaje estadístico como ruido de etiquetas dependiente de la instancia (\textit{instance-dependent label noise}) \cite{frenay_classification_2014, song_learning_2022}, provoca una divergencia severa entre la distribución de probabilidad condicional observada y la verdadera distribución física subyacente en la señal discreta:

\begin{equation}
    P(\tilde{y} = k \mid \mathbf{x}, y = c) \neq \delta_{kc}
    \label{eq:label_noise}
\end{equation}

Donde $k$ es la etiqueta nominal estática del conjunto de datos, $c$ representa la fuente electromagnética verdaderamente dominante en el subespacio temporal evaluado, $\mathbf{x} \in \mathbb{C}^N$ es el tensor fasorial de entrada, y $\delta_{kc}$ denota la delta de Kronecker.

Para cuantificar el impacto de esta asimetría probabilística en la convergencia del estimador, es necesario analizar el principio de Minimización del Riesgo Empírico (\textit{Empirical Risk Minimization}, ERM) formulado en la Teoría del Aprendizaje Estadístico \cite{vapnik_nature_1998}. Sea un estimador estadístico parametrizado por $\theta$, su objetivo teórico es minimizar el riesgo esperado $R(\theta)$ sobre la verdadera distribución conjunta subyacente $P(\mathbf{x}, y)$:

\begin{equation}
    R(\theta) = \mathbb{E}_{P(\mathbf{x}, y)} \left[ \mathcal{L}(f_\theta(\mathbf{x}), y) \right]
    \label{eq:riesgo_teorico}
\end{equation}

Donde $\mathcal{L}$ denota una función de pérdida estándar (e.g., entropía cruzada) y $f_\theta(\mathbf{x})$ es la predicción del modelo. Sin embargo, debido a la corrupción de las etiquetas en el dominio temporal extenso, el algoritmo no tiene acceso a la distribución limpia $P(\mathbf{x}, y)$, sino a una distribución empírica ruidosa $\tilde{P}(\mathbf{x}, \tilde{y})$. Por consiguiente, el optimizador minimiza un riesgo empírico corrupto $\tilde{R}(\theta)$:

\begin{equation}
    \tilde{R}(\theta) = \frac{1}{M} \sum_{i=1}^{M} \mathcal{L}(f_\theta(\mathbf{x}_i), \tilde{y}_i)
    \label{eq:riesgo_empirico}
\end{equation}

Bajo la presencia del ruido de etiquetas dependiente de la instancia formulado en la Ecuación \ref{eq:label_noise}, se demuestra teóricamente que el riesgo empírico constituye un estimador sesgado del riesgo verdadero, cumpliéndose la desigualdad fundamental $\mathbb{E}[\tilde{R}(\theta)] \neq R(\theta)$ \cite{natarajan_learning_2013, song_learning_2022}. Durante la fase de entrenamiento o ajuste paramétrico, la minimización irrestricta de la Ecuación \ref{eq:riesgo_empirico} fuerza a los parámetros $\theta$ a alinear su espacio de decisión con las características morfológicas de mayor energía presentes en la señal. Al procesar iterativamente tensores $\mathbf{x}$ gobernados por transmisiones que, en este contexto, se consideran espurias (como tramas OFDM de Wi-Fi) emparejados artificialmente con la etiqueta de un dron ($\tilde{y} = \text{Dron}$), cualquier estimador estadístico converge hacia un mínimo local igualmente espurio. El algoritmo colapsa su capacidad de generalización al asimilar los patrones de las interferencias como características discriminantes primarias de la clase objetivo, lo que se conoce en el ámbito del aprendizaje automático como \textit{feature shortcutting} o también \textit{shortcut learning}, invalidando por completo la inferencia en el límite del ruido. Esta limitación inherente al \textit{ground truth} del conjunto de datos condicionará el diseño algorítmico posterior, exigiendo la implementación de estrategias avanzadas para la mitigación de ruido de etiquetas durante el modelado.

Por consiguiente, el primer paso en la arquitectura del sistema consiste en formular un mecanismo de segmentación temporal agnóstico. Este preprocesamiento no evalúa la naturaleza del emisor, sino que se limita a rastrear la serie temporal discreta para extraer las regiones de interés (\textit{Regions of Interest}, ROI) basándose en anomalías energéticas y estadísticas. Aislar los transitorios electromagnéticos del ruido de fondo permite transformar la ventana macroscópica en tensores focalizados, paso previo y estrictamente necesario antes de abordar la clasificación y la mitigación de la corrupción de etiquetas.

\section{Detección Analítica de Ráfagas FHSS mediante Entropía Espectral de Shannon}
\label{sec:detector_entropia}

Tal y como se ha formalizado en la Sección~\ref{sec:label_noise}, la detección en el límite del ruido térmico exige un mecanismo de segmentación temporal previo al clasificador. Este preprocesador no evalúa la naturaleza del emisor, sino que rastrea la serie temporal para aislar las regiones de interés (\textit{Regions of Interest}, ROI) con concentración espectral anómala. En el contexto de los protocolos \gls{fhss}, este requisito puede satisfacerse de forma analítica explotando una propiedad fundamental de la información de Shannon \cite{zhang_entropy-based_1989}: la entropía espectral de la señal observada cae bruscamente en el instante exacto de una transmisión activa.

A diferencia de los detectores de energía convencionales, basados en la comparación de la potencia total integrada frente a un umbral fijo, el enfoque propuesto opera sobre la distribución de la energía en el dominio de la frecuencia. Este principio ha sido aplicado en trabajos previos sobre análisis espectral de señales de drones y UAVs \cite{morge-rollet_drone_2022, li_dim_2025}, y constituye la base del detector implementado en esta investigación. Los parámetros que determinan el funcionamiento de este bloque funcional no se ajustan de forma automática, ya que esto supondría una fase de entrenamiento adicional que arrastraría todos los problemas comentados en secciones anteriores, en lugar de ser una herramienta de apoyo al modelo implementado. Eneste contexto, se trata de un algoritmo determinista, cuyas constantes de operación se derivan analíticamente del estudio realizado sobre el conjunto de datos seleccionado, de las estadísticas del ruido gaussiano presente en él y de la geometría espectral de las distintas transmisiones \gls{fhss} identificadas.

\subsection{Fundamento físico: entropía como medida de concentración espectral}
\label{subsec:fundamento_entropia}

Considérese la representación tiempo-frecuencia de la señal discreta $y[n]$ obtenida mediante la \gls{stft}\footnote{La Transformada de Fourier de Tiempo Corto es una extensión de la \gls{dft}, (Transformada de Fourier en Tiempo Discreto) que permite analizar señales no estacionarias cuyas propiedades espectrales varían en el tiempo. La idea fundamental consiste en segmentar la señal en fragmentos temporales de pequeña longitud mediante una función de ventana $w[n]$ (Rectangular, Hanning, Hamming y Blackman son las más usadas \cite{zhang_window, harris_window}), aplicar la \gls{dft} a cada trama de forma independiente y organizar los espectros resultantes en una matriz bidimensional de dimensiones $\text{frecuencia} \times \text{tiempo}$. Formalmente, para una señal discreta $x[n]$, la \gls{stft} se define como $X[m, k] = \sum_{n} x[n]\, w[n - mH]\, e^{-j2\pi kn/N}$, donde $m$ es el índice de trama, $k$ el bin de frecuencia, $H$ el desplazamiento entre tramas y $N$ la longitud de la ventana de análisis. La resolución tiempo-frecuencia queda sujeta al principio de incertidumbre de Heisenberg-Gabor, ya comentado en \ref{sec:justificacion_procesamiento_IQ}: ventanas largas ofrecen alta resolución frecuencial pero baja resolución temporal, y viceversa \cite{oppenheim_discrete_2014}.}, definida en la Ecuación~\ref{eq:espectrograma}. Para cada instante temporal $m$, se dispone de un vector de densidad espectral de potencia $\mathbf{P}_m = \{P_{k,m}\}_{k=0}^{N-1}$ de longitud $N$ \textit{bins} de frecuencia. La distribución de probabilidad espectral normalizada en el \textit{frame} $m$ se define como:

\begin{equation}
    p_{k,m} = \frac{P_{k,m}}{\displaystyle\sum_{k=0}^{N-1} P_{k,m}}, \quad \forall k \in \{0, \ldots, N-1\}
    \label{eq:prob_espectral}
\end{equation}

A partir de ella, la entropía de Shannon en bits por \textit{frame} temporal se computa como:

\begin{equation}
    H(m) = -\sum_{k=0}^{N-1} p_{k,m} \cdot \log_2\!\left(p_{k,m}\right)
    \label{eq:entropia_shannon}
\end{equation}

Esta magnitud posee dos propiedades analíticas directamente explotables para la detección:

\begin{itemize}\itemsep 0em
    \item Máximo de entropía (ruido blanco): cuando la energía se distribuye de forma uniforme en todos los \textit{bins} (distribución espectral plana, característica del ruido blanco gaussiano), la distribución $\{p_{k,m}\}$ converge a la distribución uniforme y la entropía alcanza su cota superior teórica: $H_{\max} = \log_2(N)$ bits \cite{cover_elements_2006}. Para $N = 2048$ \textit{bins}, $H_{\max} = 11$ bits.

    \item Mínimo de entropía (transmisión activa): durante un salto de frecuencia \gls{fhss}, la energía se concentra en uno o pocos \textit{bins} de la portadora activa, reduciendo la distribución efectiva a un soporte mínimo. En ese instante, $H(m) \ll H_{\max}$, generando un mínimo local en la curva temporal $\{H(m)\}$ que señala inequívocamente la presencia de una ráfaga. Esta propiedad ha sido formalmente caracterizada como test de hipótesis para la detección de señales correladas en presencia de ruido blanco gaussiano aditivo \cite{zhang_entropy-based_1989}.
\end{itemize}

La idea intuitiva es directa: el ruido térmico informa en todos los \textit{bins} simultáneamente (máxima incertidumbre), mientras que el salto de frecuencia concentra la información en un subespacio reducido (mínima incertidumbre). El detector explota esta asimetría sin requerir conocimiento previo alguno sobre el protocolo, la frecuencia portadora ni la modulación empleada por el transmisor.

\subsubsection{Blanqueo espectral: la mejora clave frente a la entropía estándar}
\label{subsubsec:whitening}

La estimación directa de la entropía sobre el periodograma empírico obtenido mediante la \gls{stft} adolece de una degradación sistemática originada por la falta de planicidad espectral del canal de comunicaciones. Fenómenos como el desvanecimiento selectivo en frecuencia impuesto por el medio de propagación, las tolerancias de los osciladores y la respuesta en frecuencia no ideal de las distintas cadenas de filtros del receptor \gls{sdr} introducen alteraciones energéticas localizadas en intervalos espectrales específicos \cite{tandra_snr_2008}. Estas variaciones determinísticas inducen una entropía base inferior a la cota teórica de máxima incertidumbre $H_{\max}$, incrementando la varianza temporal del estimador y comprometiendo la detectabilidad de las ráfagas transmitidas en regímenes de baja \gls{snr}.

La solución implementada es el blanqueo espectral: antes de calcular la entropía, cada \textit{bin} se normaliza por su mediana\footnote{En estadística, la mediana muestral se prefiere frente a la media aritmética por su elevado punto de ruptura (se conoce como punto de ruptura o de quiebre de un estimador como la proporción de observaciones incorrectas que puede manejar antes de proporcionar un resultado incorrecto) del 50\,\%. En nuestro escenario, las transmisiones de alta energía (los saltos \gls{fhss} o ráfagas WiFi) actúan como valores atípicos (conocidos en el ámbito del \textit{Machine Learning} como \textit{outliers}) temporales. Con el uso de la media, estas ráfagas elevarían artificialmente el nivel de referencia del \textit{bin}, sesgando la normalización y atenuando el contraste de los saltos más débiles. La mediana permanece inalterada frente a estos eventos, estimando de forma insesgada el verdadero suelo de ruido siempre que la ocupación del canal no supere la mitad de la ventana temporal.} de potencia a lo largo del eje temporal:

\begin{equation}
    \tilde{P}_{k,m} = \frac{P_{k,m}}{\tilde{P}_{k}^{\text{med}}}, \qquad \tilde{P}_{k}^{\text{med}} = \operatorname{mediana}_{m}\!\left(P_{k,m}\right)
    \label{eq:whitening}
\end{equation}

Esta operación normaliza el ruido blanco a una potencia espectral unitaria en todos los \textit{bins}, independientemente de la forma espectral del canal. En regímenes de \gls{snr} favorables o moderadamente adversos, la aparición de un salto \gls{fhss} provoca que los \textit{bins} de la portadora activa superen de manera clara el nivel de referencia blanqueado, amplificando el contraste de la caída de entropía frente a la línea de base. Sin embargo, a medida que la señal se degrada hacia escenarios de \gls{snr} muy desfavorables (en torno a $-14$\,dB), la energía del salto queda progresivamente enmascarada por la varianza del ruido térmico, provocando que el decremento de entropía sea sutil o marginal, pero notable. El vector de entropía blanqueada, empleado en el resto del \textit{pipeline}, se calcula sustituyendo $P_{k,m}$ por $\tilde{P}_{k,m}$ en las Ecuaciones~\ref{eq:prob_espectral} y~\ref{eq:entropia_shannon}.

\subsection{Parámetros de la STFT y Resolución Tiempo-Frecuencia}
\label{subsec:parametros_stft}

La \gls{stft} se evalúa sobre la señal compleja en banda base $s[n] = y_I[n] + j\,y_Q[n]$ con ventana bilateral, preservando el espectro de frecuencias negativas necesario para representar la asimetría espectral de las señales en banda base compleja \cite{oppenheim_discrete_2014}. El parámetro de diseño fundamental es la longitud de la ventana de análisis $N_w$, que controla el equilibrio tiempo-frecuencia conforme al principio de incertidumbre de Heisenberg-Gabor:

\begin{equation}
    \Delta f = \frac{F_s}{N_w}, \qquad \Delta t \approx \frac{N_w}{2 F_s}
    \label{eq:resolucion_stft}
\end{equation}

Para el conjunto de datos \textit{NoisyUAV v2}, con $F_s = 14$\,MHz, el valor adoptado es $N_w = 2048$, que proporciona una resolución frecuencial de $\Delta f \approx 6.84$\,kHz/\textit{bin} y una resolución temporal de $\Delta t \approx 0.073$\,ms/\textit{frame}. Esta elección satisface simultáneamente dos requisitos físicos: la resolución temporal es suficiente para perfilar con precisión los flancos de transmisión de las ráfagas \gls{fhss} más cortas registradas en el conjunto de datos ($\approx 0.44$\,ms), y la resolución frecuencial permite resolver individualmente los \textit{bins} portadores de los saltos, cuyo ancho de banda en los protocolos de radiocontrol es del orden de los megahercios \cite{morge-rollet_drone_2022}.

\vspace{\baselineskip}

Para ilustrar empíricamente esta superioridad analítica, la Figura \ref{fig:ed_vs_entropy_cfar} presenta una evaluación comparativa frente a una captura del conjunto de datos operando en régimen crítico (\gls{snr}~=~$-14$\,dB). En el espectrograma superior se han identificado un total de seis transmisiones: cuatro correspondientes a saltos de frecuencia del UAV marcadas en color rojo (anotaciones 1 a 4), una ráfaga de Bluetooth en color morado (anotación 5) y una trama Wi-Fi en color negro (anotación 6). 

Al aplicar un algoritmo \gls{cfar} adaptativo moderno sobre el perfil de potencia, el detector de energía (panel central) se ve limitado por el enmascaramiento que causa el ruido térmico sumado a las interferencias: detecta únicamente las emisiones interferentes de alta potencia (Bluetooth y Wi-Fi) y omite por completo las firmas del UAV, cuya energía se encuentra por debajo del umbral del ruido de fondo. Por el contrario, el mecanismo propuesto basado en entropía espectral (panel inferior) logra segmentar con precisión tanto los cuatro saltos del dron como la ráfaga Bluetooth debido a su alta concentración frecuencial. Asimismo, ignora la emisión Wi-Fi, ya que su gran ancho de banda presenta una distribución espectral plana asimilable al ruido y no produce una caída de entropía significativa.

\begin{figure}[!htbp]
    \centering
    \includegraphics[width=\linewidth]{figuras/seccion3/ed_vs_entropy_cfar_edit.pdf}
    \caption{Comparativa entre el detector de energía y el detector de entropía espectral propuesto ante un entorno de baja relación señal a ruido y con coexistencia de interferencias de banda estrecha (Bluetooth) y banda ancha (Wi-Fi).}
    \label{fig:ed_vs_entropy_cfar}
\end{figure}
